\chapter*{LEMBAR PENGESAHAN}
\addcontentsline{toc}{chapter}{LEMBAR PENGESAHAN}

\begin{singlespacing}

\begin{center}
  {\large\bfseries
    Pemanfaatan Generative AI untuk Process Discovery pada Data
Operasional Pemeliharaan Perusahaan Listrik
  \par}

  \vspace{0.3cm}

   {\large\bfseries
   Utilization of Generative AI for Process Discovery in
Operational Maintenance Data Electric Power Company
   \par}
\end{center}

\vspace{0.3cm}

\begin{center}
  {\large\bfseries 103012480032\par}
  {\large\bfseries Sakha Wibisono\par}
\end{center}

\vspace{0.3cm}

\begin{center}
  Tugas akhir ini telah diterima dan disahkan untuk memenuhi sebagai syarat memperoleh
  gelar pada Program Studi Sarjana Informatika\\
  Fakultas Informatika\\
  \vspace{0.2cm}
  Universitas Telkom
\end{center}

\vspace{0.3cm}

\begin{center}
  Bandung, 07 April 2026
\end{center}

\vspace{0.2cm}
\begin{center}
  Menyetujui,
\end{center}

\vspace{1.0cm}

\begin{center}
  Pembimbing I, \\[2.0cm]
  \textbf{Angelina Prima Kurniati, S.T., M.T., Ph.D.}\\
  NIP: 06830027
\end{center}

\vspace{1.0cm}

\begin{center}
  Ketua Program Studi Sarjana Informatika\\
  Fakultas Informatika,\\[2.0cm]
  \textbf{MAHMUD DWI SULISTIYO, S.T., M.T., Ph.D.}\\
  NIP: 13880017
\end{center}

\end{singlespacing}
